\subsubsection{Hamiltonian path / cycle (bitmask DP + backtracking)}
\begin{lstlisting}[style=mystyle, language=C++]
// TC: O(N^2 * 2^N)  (bitmask DP), O(N!) peor caso (backtracking)
// usa: existencia y construccion de camino / ciclo hamiltoniano;
//     TSP (min costo de ciclo hamiltoniano) si las aristas tienen peso
#include <bits/stdc++.h>
using namespace std;

// -- Teorema de Dirac (condicion SUFICIENTE) --
//   Si G es simple, |V| = n >= 3, y todo vertice tiene grado >= n/2,
//   entonces G tiene un ciclo hamiltoniano. (No es necesario.)
//   Ore: si para todo par no adyacente (u,v) se cumple
//        deg(u) + deg(v) >= n, entonces G es hamiltoniano.

bool satisfiesDirac(int n, const vector<vector<int>>& adj){
	if(n < 3) return false;
	for(int v=0;v<n;v++) if((int)adj[v].size() < (n+1)/2) return false;
	return true;
}
bool satisfiesOre(int n, const vector<vector<int>>& adj){
	if(n < 3) return false;
	// construye conjunto de adyacencia
	vector<vector<char>> A(n, vector<char>(n, 0));
	for(int u=0;u<n;u++) for(int v:adj[u]) A[u][v] = 1;
	for(int u=0;u<n;u++)
		for(int v=u+1;v<n;v++)
			if(!A[u][v] && ((int)adj[u].size() + (int)adj[v].size() < n))
				return false;
	return true;
}

// -- Bitmask DP: Hamiltonian path starting at 'src' --
// dp[mask][v] = true si hay un camino que visita exactamente 'mask' y termina en v
// memoria: O(N * 2^N) bool; tiempo: O(N^2 * 2^N)
bool hasHamiltonianPathFrom(int n, const vector<vector<int>>& adj, int src){
	vector<vector<char>> dp(1<<n, vector<char>(n, 0));
	dp[1<<src][src] = 1;
	for(int mask=1; mask<(1<<n); mask++){
		for(int v=0; v<n; v++) if(dp[mask][v]){
			for(int u : adj[v]){
				if(!(mask & (1<<u))){
					dp[mask | (1<<u)][u] = 1;
				}
			}
		}
	}
	for(int mask=1; mask<(1<<n); mask++){
		if(__builtin_popcount(mask) != n) continue;
		for(int v=0; v<n; v++) if(dp[mask][v]) return true;
	}
	return false;
}

// -- Bitmask DP: Hamiltonian cycle --
// dp[mask][v] = true si hay un camino de 'src' que visita 'mask' y termina en v
bool hasHamiltonianCycle(int n, const vector<vector<int>>& adj, int src = 0){
	if(n < 3) return false;
	vector<vector<char>> dp(1<<n, vector<char>(n, 0));
	dp[1<<src][src] = 1;
	for(int mask=1; mask<(1<<n); mask++){
		for(int v=0; v<n; v++) if(dp[mask][v]){
			for(int u : adj[v]){
				if(!(mask & (1<<u))){
					dp[mask | (1<<u)][u] = 1;
				}
			}
		}
	}
	int full = (1<<n) - 1;
	for(int v=0; v<n; v++){
		if(dp[full][v]){
			// cierra el ciclo si v es adyacente a src
			for(int u : adj[v]) if(u == src) return true;
		}
	}
	return false;
}

// -- Reconstruccion de un camino hamiltoniano (bitmask DP) --
// devuelve la secuencia de vertices o vacio si no existe
vector<int> findHamiltonianPath(int n, const vector<vector<int>>& adj, int src){
	vector<vector<int>> par(1<<n, vector<int>(n, -1));
	vector<vector<char>> dp(1<<n, vector<char>(n, 0));
	dp[1<<src][src] = 1;
	for(int mask=1; mask<(1<<n); mask++){
		for(int v=0; v<n; v++) if(dp[mask][v]){
			for(int u : adj[v]){
				if(!(mask & (1<<u))){
					dp[mask | (1<<u)][u] = 1;
					par[mask | (1<<u)][u] = v;
				}
			}
		}
	}
	int endV = -1, full = -1;
	for(int mask=1; mask<(1<<n); mask++){
		if(__builtin_popcount(mask) != n) continue;
		for(int v=0; v<n; v++) if(dp[mask][v]){ endV = v; full = mask; break; }
		if(endV != -1) break;
	}
	if(endV == -1) return {};
	vector<int> path;
	int cur = endV, m = full;
	while(cur != -1){
		path.push_back(cur);
		int p = par[m][cur];
		m &= ~(1<<cur);
		cur = p;
	}
	reverse(path.begin(), path.end());
	return path;
}

// -- TSP minimo: ciclo hamiltoniano de costo minimo --
// dp[mask][v] = min costo de camino src -> ... -> v visitando 'mask'
// respuesta = min_v dp[full][v] + w(v, src) si el cierre es valido
const long long LINF = (1LL<<60);
long long tspMinCycle(int n, const vector<vector<long long>>& w, int src = 0){
	if(n == 1) return 0;
	vector<vector<long long>> dp(1<<n, vector<long long>(n, LINF));
	dp[1<<src][src] = 0;
	for(int mask=1; mask<(1<<n); mask++){
		for(int v=0; v<n; v++){
			if(dp[mask][v] == LINF) continue;
			for(int u=0; u<n; u++){
				if(!(mask & (1<<u)) && w[v][u] < LINF){
					dp[mask | (1<<u)][u] = min(dp[mask | (1<<u)][u],
					dp[mask][v] + w[v][u]);
				}
			}
		}
	}
	int full = (1<<n) - 1;
	long long ans = LINF;
	for(int v=0; v<n; v++){
		if(dp[full][v] < LINF && w[v][src] < LINF){
			ans = min(ans, dp[full][v] + w[v][src]);
		}
	}
	return ans;
}

int main(){
	// ejemplo: ciclo de 4 vertices (0-1-2-3-0) -> claramente hamiltoniano
	int n = 4;
	vector<vector<int>> adj(n);
	auto add = [&](int u, int v){
		adj[u].push_back(v); adj[v].push_back(u);
	};
	add(0,1); add(1,2); add(2,3); add(3,0);

	cout << "Dirac: " << satisfiesDirac(n, adj) << " (n/2 = 2, todos grado 2 -> true)\n";
	cout << "Ore:   " << satisfiesOre(n, adj) << "\n";
	cout << "cycle? " << hasHamiltonianCycle(n, adj) << "\n";
	cout << "path?  " << hasHamiltonianPathFrom(n, adj, 0) << "\n";

	auto path = findHamiltonianPath(n, adj, 0);
	for(int v : path) cout << v << " ";
	cout << "\n";

	// TSP: mismo grafo con pesos -> ciclo minimo
	vector<vector<long long>> w(n, vector<long long>(n, LINF));
	for(int i=0;i<n;i++) w[i][i] = 0;
	w[0][1] = w[1][0] = 10;
	w[1][2] = w[2][1] = 20;
	w[2][3] = w[3][2] = 30;
	w[3][0] = w[0][3] = 40;
	cout << "TSP min = " << tspMinCycle(n, w) << "\n"; // 10+20+30+40 = 100
	return 0;
}
\end{lstlisting}
